\section{Conclusion}

In this work, we addressed the critical challenge of predicting ADMET properties in the extreme low-data regime ($P \gg N$) characteristic of modern computational cheminformatics. By evaluating models on the Caco-2 cell permeability benchmark utilizing dense, 768-dimensional MoLFormer representations, we highlighted the necessity of moving beyond mere point predictions toward mathematically rigorous epistemic uncertainty quantification for Out-of-Distribution (OOD) detection.

Our extensive Optuna benchmarking revealed fundamental limitations in classical tree-based methodologies when applied to deep representation learning. Despite massive hyperparameter optimization and Deep Ensembling to artificially induce variance, XGBoost suffered from a severe \textit{geometric mismatch}. The reliance on orthogonal, axis-aligned splits to partition a continuous, distributed latent space led to fragmented learning, rapid overfitting, and a catastrophic collapse in predictive variance. Consequently, the Deep Ensemble yielded pathologically uncalibrated prediction intervals (PICP\_95 $\approx$ 21\%) and failed entirely to rank its own errors.

Conversely, we empirically demonstrated that Stochastic Variational Deep Kernel Learning (SV-DKL) is theoretically and practically superior for parsing foundation model embeddings. By natively computing spatial covariance across the entire high-dimensional vector simultaneously, the Gaussian Process effectively interpolates complex molecular geometries without succumbing to the curse of dimensionality. 

Furthermore, our ablation study conclusively proved that applying a bi-Lipschitz constraint via Spectral Normalization is strictly necessary to prevent \textit{feature collapse}. While unregularized deep networks blindly squash OOD samples into the training manifold—resulting in overconfident errors and flat rejection curves—Spectral Normalization successfully restores distance awareness. This allows the SV-DKL model to accurately recognize unfamiliar molecules, correctly inflating its predictive variance and drastically improving both Negative Log-Likelihood (NLL) and the Area Under the Rejection Curve (AURC).

Ultimately, this study establishes that combining expressive foundation model representations with properly regularized probabilistic architectures bridges the gap between predictive power and safety. SV-DKL with Spectral Normalization not only achieves superior point accuracy over optimized classical ensembles, but it provides the calibrated, highly discriminative epistemic uncertainty required for safe, real-world deployment in high-stakes pharmaceutical drug discovery.